\chapter{Conclusion}
\label{chap:conclusion}

This thesis addressed the problem of deploying a standards-compliant sensor data
management platform for Czech University of Life Sciences Prague within the European
eLTER research infrastructure. The existing UFZ Time.IO solution provided per-Thing
schema isolation, MQTT-driven worker pipelines, and OGC~STA data access, but lacked a
usable sensor management interface, project-level access control, and a reproducible
deployment mechanism suitable for external adopters. This work designed, implemented,
and validated a three-layer extension that fills these gaps.

\section*{Summary of Results}
\addcontentsline{toc}{section}{Summary of Results}

The implemented system consists of three independently deployable components that together
form a complete environmental monitoring platform:

\begin{enumerate}
  \item \textbf{TSM Orchestration Adaptation.}
        The upstream \texttt{tsm-orchestration} stack was adapted with local STA views
        implementing the missing FROST compatibility layer, a modular Docker~Compose
        structure enabling overlay-based configuration, and Podman rootless support for
        security-constrained infrastructure. Eighteen marked workarounds bypass the
        non-functional upstream SMS service.

  \item \textbf{The Water Data Platform.}
        A new application layer built with FastAPI (Python) and Next.js~15 (TypeScript)
        provides project-oriented sensor management, role-based access control with
        Keycloak integration, MQTT-driven alert evaluation, automated QA/QC via SaQC,
        geospatial layer integration via GeoServer, and a multilingual web portal in
        English, Czech, and Slovak.

  \item \textbf{The Unified Deployment Layer.}
        A deployment orchestrator that merges the two stacks via Docker~Compose
        multi-file merge, manages secrets, validates cross-stack configuration
        consistency, and provides a scripted operator interface via a Makefile.
        A single \texttt{make up} command starts the full 22-container stack.
\end{enumerate}

\noindent
The platform was validated through three complementary approaches:

\begin{itemize}
  \item \textbf{Automated testing.} A backend suite of 804~unit and integration tests
        covers the service layer, API endpoints, middleware, and schema validation
        (Section~\ref{sec:test-approach-backend}). All 20~functional requirements are
        fully met (Section~\ref{sec:test-requirements}).

  \item \textbf{Load testing.} Locust-based load tests with 10~and 25~concurrent users
        confirmed that all API endpoints respond within the 2-second threshold required
        by NFR-5. Lightweight CRUD operations maintain sub-15\,ms median latency;
        data-heavy endpoints involving FROST-Server fan-out remain under 310\,ms at the
        95th~percentile. The system sustains 11.7~requests/second under 25~concurrent
        users (Section~\ref{sec:test-performance}).

  \item \textbf{Deployment validation.} The full stack was deployed on CZU
        infrastructure and verified using the automated \texttt{check.sh} script, which
        confirmed container health, endpoint reachability, and cross-service
        connectivity for all 22~containers
        (Section~\ref{sec:test-deployment-czu}).
\end{itemize}

The design choice that left the deepest mark on the result was keeping
the ingestion layer and the data platform as separately deployable stacks instead of
merging them into one monolithic codebase. That separation made it possible to pull in
upstream TSM updates while letting the application layer evolve on its own. In practice,
however, the upstream Time.IO project proved unresponsive to external adopters:
the SMS service stayed broken, the frontends emitted invalid MQTT messages, and integration
documentation was absent. The data platform exists precisely because of these shortcomings.

\section*{Limitations}
\addcontentsline{toc}{section}{Limitations}

The biggest limitation is the continued reliance on \texttt{TSM-TEMP}
workarounds. Across \texttt{tsm-orchestration}, 18~marked locations bypass the non-functional
upstream SMS service with local PostgreSQL views. These views work well operationally
(Section~\ref{sec:impl-tsm-assessment}), but they carry the label of temporary fixes against a
codebase whose upstream shows no sign of addressing the root cause.

Backend test coverage, though extensive for the core service
layer (797~tests; Section~\ref{sec:test-approach-backend}), has gaps in the FROST client integration
(14\% coverage for \texttt{async\_frost\_client}) and the SMS service adapter (14\%). These
modules depend on external services that are hard to exercise without a running
FROST-Server and Keycloak instance.

Performance has not been benchmarked at scale beyond the CZU sensor fleet. The primary
bottleneck identified, HTTP request fan-out to FROST-Server's per-thing context paths
(Section~\ref{sec:test-perf-bottleneck}), would become critical at deployments with
thousands of Things. Quantitative throughput measurements under realistic multi-site load
remain future work.

Finally, the platform is not yet eLTER-ready. The architectural foundations are aligned
(OGC~STA data access, OIDC authentication, containerized deployment), but the
eLTER-specific requirements identified in Section~\ref{sec:test-deployment-elter} ---
metadata schemas, cross-site federation, vocabulary alignment, and FAIR compliance ---
represent substantial development effort beyond the scope of this thesis.

\section*{Future Work}
\addcontentsline{toc}{section}{Future Work}

What this thesis points toward is a clear next step: instead
of continuing to patch the upstream Time.IO codebase through workarounds, the more
sustainable path is to extract the \emph{architectural ideas} of Time.IO --- per-Thing schema
isolation, MQTT-driven worker pipelines, OGC~STA as the interoperability standard --- and
rebuild the problematic components from scratch, purpose-built for eLTER adoption.

The concrete future work items are:

\begin{enumerate}
  \item \textbf{TSM component rewrite.} Replace the \texttt{TSM-TEMP} workarounds
        and the non-functional SMS service with a clean STA view layer that is maintained
        as part of the platform rather than labeled as temporary. The rewrite scope is
        well-defined: 18~marked locations in one Python module and one SQL directory.

  \item \textbf{eLTER metadata integration.} Extend the project and sensor models
        with DEIMS-SDR site identifiers and EnvThes vocabulary bindings so that
        observations link to the eLTER site registry and use semantically aligned
        property descriptions. Implement metadata export conforming to the eLTER
        dataset model so that datasets can be registered in the Digital Asset
        Register and the STA endpoint listed in the Online Asset Register.

  \item \textbf{eLTER~LogIn federation.} Configure the local Keycloak instance as
        an identity broker for the eLTER~LogIn service (Perun~AAI), mapping eLTER
        roles to the platform's role model to enable single-sign-on for eLTER users.

  \item \textbf{Cross-site federation.} Implement a federation layer for aggregated
        queries across independently operated platform instances, a prerequisite for
        eLTER's distributed, multi-site architecture.

  \item \textbf{FAIR data compliance.} Add persistent identifiers (DOIs), structured
        provenance records, and machine-readable licensing metadata to support the FAIR
        data principles mandated by eLTER.

  \item \textbf{Performance optimization.} Introduce a caching layer (Redis) for
        frequently accessed FROST metadata, evaluate FROST-Server alternatives or
        clustering, and conduct load testing at scale beyond the CZU sensor fleet.

  \item \textbf{Test coverage completion.} Extend the automated test suite to cover the
        FROST client integration and SMS service adapter, potentially using
        containerized test fixtures (Testcontainers) for FROST-Server and Keycloak.
\end{enumerate}

\medskip\noindent
The OGC SensorThings API proved to be a well-chosen interoperability standard for this class
of environmental monitoring platform. Its entity model maps naturally onto the sensor data
management domain, and the availability of the FROST-Server implementation reduced the
initial integration effort. The primary friction was not with the standard itself but with
the state of the upstream implementation layer that was intended to bridge it to operational
sensor networks. The experience of this thesis suggests that the concept behind Time.IO ---
a modular, MQTT-driven, STA-compliant sensor data platform --- is sound, but a purpose-built
implementation targeting eLTER's specific requirements would be more effective than continued
adaptation of an upstream codebase that does not prioritise external adopters.
